% 原书第233--261页；第七篇及L. Illusie附录。
\chapter{$\ell$进上同调的有限性定理}\label{VII}
\origpage{239}{233}
\section{定理的陈述}\label{VII:1}
本篇始终设 $S$ 为\term{诺特概形}{Noetherian scheme}{schéma noethérien}，$A$ 为左诺特环，且被某个在 $S$ 上可逆的整数零化。我们的主要结果如下。
\begin{theorem_}\label{VII:1-1}
设 $S$ 正则，维数为 $0$ 或 $1$。设 $f:X\to Y$ 是有限型 $S$概形之间的态射，$\sF$ 是 $X$ 上的可构造左 $A$模层。则各层 $\eR^i f_\ast\sF$ 都可构造。
\end{theorem_}
证明见第\ref{VII:2}节及\ref{VII:3-10}。
\subsection{注}\label{VII:1-2}
在特征 $0$ 时，这个结果不如\cite[XIX，第5节]{sga4}一般；该处对特征 $0$ 的优良概形之间任意有限型态射证明了上述结论。\cite[XIX]{sga4}的证明一方面使用奇点消解，另一方面使用概形等特征这一条件，以便从消解推出“纯性定理”。
\subsection{注}\label{VII:1-3}
若复形 $K\in\ob\eD(X,A)$ 的上同调层都可构造，就称它\term{可构造}{constructible}{constructible}；以一个下标 $c$ 表示 $\eD(X,A)$（或 $\eD^+$、$\eD^-$、$\eD^b$）中由可构造复形组成的子范畴。下文将自由使用导出范畴的语言；在这一语言中，\ref{VII:1-1}表示 $\eR f_\ast:\eD^+(X,A)\to\eD^+(Y,A)$ 将 $\eD_c^+$ 映入 $\eD_c^+$。具体而言：
\begin{enumerate}[\indent a)]
\item 当 $K$ 仅在 $0$ 次有层 $\sF$ 时，$\sH^i\eR f_\ast K=\eR^i f_\ast\sF$，故导出范畴中的陈述蕴含定理。
\item 反过来，使用谱序列
\[
 E_1^{p q}=\eR^q f_\ast\sH^p(K)\Rightarrow\sH^{p+q}\eR f_\ast K.
\]
\end{enumerate}
导出范畴语言的优点是，可以用简单的传递性公式 $\eR(fg)_\ast=\eR f_\ast\eR g_\ast$ 代替 Leray 谱序列
$E_2^{pq}=\eR^p f_\ast\eR^q g_\ast\Rightarrow\eR^{p+q}(fg)_\ast$。
在\ref{VII:1-1}的假设下，由\cite[X]{sga4}知 $\eR f_\ast$ 的\term{上同调维数}{cohomological dimension}{dimension cohomologique}有限。因此，上面可以将 $\eD^+$ 换成 $\eD$、$\eD^-$ 或 $\eD^b$。
\subsection{证明的思路}\label{VII:1-4}
\begin{enumerate}[\indent a)]
\item Poincaré 对偶处理 $X$ 光滑、$\sF$ 局部常值且 $Y=S$ 的情形。维数假设用在 $S$ 上的 $\rHom$ 计算中。
\item 将 $f$ 分解成 $gj$，其中 $g$ 固有，$j$ 为开嵌入。于是 $\eR f_\ast=\eR g_\ast\eR j_\ast$，固有态射的有限性定理控制 $\eR g_\ast$。经过分解归约，可设 $X$ 光滑、$\sF$ 局部常值、$f$ 为开嵌入且 $Y$ 在 $S$ 上固有。
\item 对 $\dim X$ 归纳，同时允许改变 $S$，大致可假定定理在 $Y$ 的某个于 $S$ 上有限的部分之外成立。记结构态射为 $a:X\to S$、$b:Y\to S$；这时证明：若对 $\sF$ 和 $\eR f_\ast$ 定理不成立，则对 $\sF$ 和 $\eR b_\ast\eR f_\ast=\eR a_\ast$ 也不成立，矛盾。
\end{enumerate}
\begin{corollary_}\label{VII:1-5}
在定理的假设下，四种运算 $\eR f_\ast$、$\eR f_!$、$f^\ast$、$\eR f^!$ 将范畴 $\eD_c(X,A)$ 和 $\eD_c(Y,A)$ 中的对象沿相应方向映入另一个范畴。
\end{corollary_}
$\eR f_\ast$ 已在上面处理，$\eR f_!$ 见\cite[XVII 5.3]{sga4}，$f^\ast$ 显然。只剩 $\eR f^!$。问题是局部的，所以可设 $f$ 分解为 $X\xrightarrow i Z\xrightarrow gY$，其中 $i$ 为闭嵌入，$g$ 光滑且纯相对维数为 $n$。Poincaré 对偶公式
$\eR g^!K=g^\ast K(n)[2n]$ 与传递性 $\eR f^!=\eR i^!\eR g^!$，将问题归结为对 $i$ 证明\ref{VII:1-5}。
对任意 $L\in\eD(Z,A)$，若 $j:U=Z\setminus X\hookrightarrow Z$ 为包含映射，则有三角
\[
 i_\ast\eR i^!L\longrightarrow L\longrightarrow\eR j_\ast j^\ast L\longrightarrow(i_\ast\eR i^!L)[1].
\]
因而对 $i$ 的\ref{VII:1-5}由对 $j$ 的\ref{VII:1-1}推出。\translatornote{原书此处的 Poincaré 对偶公式有漏字，并将 $L$ 写成了 $K$；这里展开了所用的局部化三角。若以通常的映射锥 $\operatorname{Cone}(u)$ 为约定，则 $i_\ast\eR i^!L=\operatorname{Cone}(L\to\eR j_\ast j^\ast L)[-1]$；这一移位不影响可构造性的论证。}
\begin{corollary_}\label{VII:1-6}
设 $X$ 如定理所述，且 $A$ 交换。若 $\sF$、$\sG$ 是 $X$ 上可构造的 $A$模层，则各 $\Ext_A^i(\sF,\sG)$ 都可构造；亦即 $\rHom$ 将 $\eD_c^-(X,A)\times\eD_c^+(X,A)$ 映入 $\eD_c^+(X,A)$。
\end{corollary_}
对 $\sF$ 分解归约，可设 $\sF=j_!\sF_1$，其中 $j:Y\hookrightarrow X$ 是局部闭嵌入，$\sF_1$ 在 $Y$ 上局部常值。此时
\[
 \rHom(j_!\sF_1,\sG)=\eR j_\ast\rHom(\sF_1,\eR j^!\sG).
\]
根据\ref{VII:1-3}，归结为 $\sF$ 局部常值的情形；问题既是局部的，还可设其常值。若 $\sF$ 局部常值可构造，则 $\Hom(\sF,\sG)_x=\hom(\sF_x,\sG_x)$，导出内部 Hom 也有相同性质。因此，$\sF$ 常值时，可用其常值模的一个每项有限生成的投射分解计算各 $\Ext^i$；当 $\sG$ 可构造时，这些 $\Ext^i$ 也可构造，推论得证。\translatornote{原式右端的内部 Hom 应取导出函子；否则对非投射的常值模不成立。例如 $A=\mathbf Z/4$，在一点上取 $\sF_1=\sG=\mathbf Z/2$，则 $\operatorname{Ext}^1_A(\mathbf Z/2,\mathbf Z/2)\ne0$，而普通 Hom 集中在 $0$ 次。上式采用导出内部 Hom 后，紧接着的投射分解论证即适用。}
\subsection{注}\label{VII:1-7}
因为 $\eR f_\ast$ 的上同调维数有限，它将 Tor 维数有限（分别为 $\leqslant d$）的复形映为 Tor 维数有限（分别为 $\leqslant d$）的复形，见\cite[XVII 5.2.11]{sga4}。由\cite[XVII 5.2.10]{sga4}和\ref{VII:1-5}的证明知，有限 Tor 维数在四种运算 $\eR f_\ast$、$\eR f_!$、$f^\ast$、$\eR f^!$ 下稳定。对\ref{VII:1-6}的论证略作变形：沿 $X$ 的分层对 $K$ 分解归约，使其上同调层局部常值，再局部化，将 $K$ 换成一个由有限型局部自由层组成的有限复形。于是，当 $A$ 交换时，$\rHom$ 诱导
\[
 \rHom:\eD_{\mathrm{ctf}}^b(X,A)\times\eD_{\mathrm{tf}}^b(X,A)\longrightarrow\eD_{\mathrm{tf}}^+(X,A),
\]
这里 t.f. 表示\term{有限Tor维数}{finite Tor dimension}{tor-dimension finie}。
\subsection{}\label{VII:1-8}
在定理的假设下，同一方法还给出\term{局部双重对偶定理}{local biduality theorem}{théorème de bidualité locale} $K\iso DDK$（第\ref{VII:4}节）。我们也将证明附近循环层的有限性定理。对一般的 $S$，仍有如下“泛性”定理。
\begin{theorem_}\label{VII:1-9}
设 $f:X\to Y$ 是有限型 $S$概形之间的态射，$\sF$ 是 $X$ 上的可构造 $A$模层。存在 $S$ 的稠密开集 $U$，使得：
\begin{enumerate}[\indent (i)]
\item 在 $U$ 上方，各 $\eR^i f_\ast\sF$ 可构造，且除有限多个外都为零。
\item $\eR^i f_\ast\sF$ 的形成与任意基变换 $S'\to U\subset S$ 相容。
\end{enumerate}
\end{theorem_}
若 $S$ 是域的谱，则 $U=S$。
当 $S$ 为域的谱时，通过取极限可知：对 $S$概形之间任意拟紧拟分离态射及任意层 $\sF$，$\eR^i f_\ast$ 与 $S$上的基变换相容。
\begin{corollary_}\label{VII:1-10}
设 $X$ 在可分闭域 $k$ 上有限型，$\sF$ 为可构造 $A$模层，则各 $\h^i(X,\sF)$ 都是有限生成模。
\end{corollary_}
这是 $S=\spec(k)$、$Y=S$ 的特例。
\begin{corollary_}\label{VII:1-11}
设 $X$、$Y$ 为可分闭域 $k$ 上的有限型概形，$K\in\ob\eD^-(X,A^\circ)$、$L\in\ob\eD^-(Y,A)$ 分别为右模层复形和左模层复形，则 Künneth 态射
\[
 \eR\Gamma(X,K)\lotimes\eR\Gamma(Y,L)\longrightarrow\eR\Gamma(X\times Y,\pr_1^\ast K\lotimes\pr_2^\ast L)
\]
是同构。
\end{corollary_}
证明与\cite[XVII 5.4.3]{sga4}相同。考虑
\[\xymatrix{X\times Y\ar[r]^-{\pr_1}\ar[d]^-{\pr_2}&X\ar[d]^-a\\Y\ar[r]^-b&\spec(k).}\]
沿 $b$ 作基变换，得 $\eR\pr_{2,\ast}\pr_1^\ast K=b^\ast\eR\Gamma(X,K)$。于是，由\cite[XVII 5.2.11]{sga4}及其证明，
\begin{align*}
 \eR\Gamma(X\times Y,\pr_1^\ast K\lotimes\pr_2^\ast L)
 &=\eR\Gamma\bigl(Y,\eR\pr_{2,\ast}(\pr_1^\ast K\lotimes\pr_2^\ast L)\bigr)\\
 &=\eR\Gamma\bigl(Y,(\eR\pr_{2,\ast}\pr_1^\ast K)\lotimes L\bigr)\\
 &=\eR\Gamma\bigl(Y,b^\ast\eR\Gamma(X,K)\lotimes L\bigr)\\
 &=\eR\Gamma(X,K)\lotimes\eR\Gamma(Y,L).
\end{align*}
\subsection{}\label{VII:1-12}
最后，作为\ref{VII:1-9}(ii)的局部对应，我们将证明泛局部无环性定理（\ref{VII:2-13}）。

\origpage{244}{238}
\section{泛性定理}\label{VII:2}
本节证明\ref{VII:1-9}及泛局部无环性定理\ref{VII:2-13}。证明\ref{VII:1-9}时，立即可归结为 $S$ 整的情形。记其泛点为 $\eta$。
\subsection{}\label{VII:2-1}
先在以下附加假设下证明\ref{VII:1-9}：$X$ 在 $S$ 上光滑，纯相对维数为 $n$，$A=\dZ/m$，$\sF$ 局部常值，且 $Y=S$。令 $\sF'=\Hom(\sF,\dZ/m)$。我们证明：若各 $\eR^i f_!\sF'$ 局部常值，则\ref{VII:1-9}的结论对 $U=S$ 成立。回忆，若 $\sL$ 局部常值可构造，则各 $\Ext^i(\sL,\sG)$ 可逐茎计算，且当 $\sG$ 可构造时也可构造。还要用到 $\dZ/m$ 是内射 $\dZ/m$模，并且是对偶化模，故 $\sF=\rHom(\sF',\dZ/m)$。Poincaré 对偶
\[
 \eR f_\ast\rHom(K,\eR f^!L)=\rHom(\eR f_!K,L)
\]
取 $K=\sF'$、$L=\dZ/m$，并利用 $\eR f^!L=\dZ/m[2n](n)$，给出
\[
 \eR^{2n-i} f_\ast\sF=\Hom(\eR^i f_!\sF',\dZ/m)(-n).
\]
右边的层局部常值，且其形成与任意基变换相容，断言即得。
\subsection{}\label{VII:2-2}
再在 $X$ 在 $S$ 上光滑、$\sF$ 局部常值、$Y=S$ 的附加假设下证明\ref{VII:1-9}。
\begin{enumerate}[a)]
\item 将 $X$ 分成连通分支，可设其在 $S$ 上纯相对维数为 $n$。
\item 将 $A$ 分成乘积，可设 $\ell^mA=0$，其中素数 $\ell$ 在 $S$ 上可逆。再用各 $\ell^k\sF$ 滤过 $\sF$，借助对应谱序列归结为 $\ell\sF=0$。此时可将 $A$ 换成 $A/\ell A$，即设 $\ell A=0$。
\item 将 $S$ 换成适当的 $U$，可设存在有限平展 Galois 覆盖 $X_1/X$，其 Galois 群为 $G$，使 $\sF$ 在 $X_1$ 上的拉回 $\sF_1$ 为常层，常值为 $F$。记 $f_1:X_1\to S$ 为投影，则
\[
 \eR^i {f_1}_\ast\sF_1=(\eR^i {f_1}_\ast\dZ/\ell)\otimes_{\dZ/\ell}F.
\]
还有 Hochschild--Serre 谱序列，也就是覆盖 $X_1/X$ 的 Leray 谱序列：
\[
 \underline\h^p(G,\eR^q {f_1}_\ast\sF_1)\Rightarrow\eR^{p+q}f_\ast\sF.
\]
根据\ref{VII:2-1}，缩小 $U$ 后，可设 $\eR^i {f_1}_\ast\dZ/\ell$ 局部常值，且其形成与任意基变换相容。各 $\eR^i f_\ast\sF$ 因而也有这一性质。最后，若 $\bar\eta$ 为位于 $S$ 的泛点 $\eta$ 上的几何点，则各 $\eR^i f_\ast\sF$ 除有限多个外均为零，因为其茎 $(\eR^i f_\ast\sF)_{\bar\eta}=\h^i(X_{\bar\eta},\sF)$ 具有这一性质。
\end{enumerate}
\subsection{}\label{VII:2-3}
对 $n$ 归纳证明
\begin{equation*}\tag{$\ast_n$}\label{VII:eq:2-3-1}
 \begin{gathered}
 \text{当 $\dim X_\eta\leqslant n$，且 $f$ 为像稠密的开嵌入时，}\\
 \text{\ref{VII:1-9}的结论成立。}
 \end{gathered}
\end{equation*}
当 $n=0$ 时，缩小 $S$ 后有 $X=Y$，所以 $(\ast_0)$ 显然。假设 $(\ast_{n-1})$，来证明\eqref{VII:eq:2-3-1}。在 $(\ast_{n-1})$ 中，可将“开嵌入”换成“嵌入”：只需分解为开嵌入和闭嵌入即可。
\begin{lemma_}\label{VII:2-4}
缩小 $S$ 后，\ref{VII:1-9}的结论在某个 $Y'\subset Y$ 上方成立，且 $Y'$ 的补集 $Y_1$ 在 $S$ 上有限。
\end{lemma_}
断言对 $Y$ 是局部的，可设 $Y$ 仿射，嵌入 $Y\subset\dA_S^N$。将归纳假设 $(\ast_{n-1})$ 用于
\[\xymatrix{X\ar[r]^-f\ar[dr]&Y\ar[d]^-{\pr_i}\\&\dA_S^1,}\]
便知对每个 $i$，存在 $\dA_S^1$ 的稠密开集 $U_i$，使\ref{VII:1-9}的结论在 $\pr_i^{-1}(U_i)$ 上方成立；于是它在各 $\pr_i^{-1}(U_i)$ 的并上方也成立，得到\ref{VII:2-4}。
\addtocounter{subsection}{1}
\subsection{}\label{VII:2-6}
先在 $X$ 在 $S$ 上光滑且 $\sF$ 在 $X$ 上局部常值的情形证明\eqref{VII:eq:2-3-1}。问题对 $Y$ 是局部的，因此可先设 $Y$ 仿射，再将其换成在射影空间中的闭包而设其射影。取\ref{VII:2-4}给出的 $i:Y_1\hookrightarrow Y$ 和 $j:Y'\to Y$：
\[\xymatrix{X\ar@{^{(}->}[r]\ar[dr]_-a&Y\ar[d]^-b&Y_1\ar[l]_-i\ar[dl]^-{b_1}\\&S.}\]
缩小 $S$ 后，已知 $j^\ast\eR f_\ast\sF$ 可构造，且其形成与 $S$ 上的任意基变换相容；也知道 $\eR a_\ast\sF=\eR b_\ast\eR f_\ast\sF$ 具有同样性质。
对短正合列
\begin{equation*}\tag{1}\label{VII:eq:2-6-1}
 0\longrightarrow j_!j^\ast\eR f_\ast\sF\longrightarrow\eR f_\ast\sF\longrightarrow i_!i^\ast\eR f_\ast\sF\longrightarrow0
\end{equation*}
所定义的三角施以 $\eR b_\ast$，得三角
\begin{equation*}\tag{2}\label{VII:eq:2-6-2}
 \cdots\longrightarrow\eR b_\ast j_!j^\ast\eR f_\ast\sF\longrightarrow\eR a_\ast\sF\longrightarrow {b_1}_\ast i^\ast\eR f_\ast\sF\longrightarrow\cdots.
\end{equation*}
前两项可构造，且其形成与 $S$ 上任意基变换相容；第一项使用了固有态射 $b$ 的有限性定理。因此第三项也如此。由于 $b_1$ 有限，可知 $i^\ast\eR f_\ast\sF$ 可构造且其形成与 $S$ 上任意基变换相容；再由\eqref{VII:eq:2-6-1}知 $\eR f_\ast\sF$ 亦然。
\subsection{}\label{VII:2-7}
现在证明一般情形的\eqref{VII:eq:2-3-1}。先归结为 $X$ 中存在在 $S$ 上光滑的稠密开集 $V$ 的情形。若 $S$ 为完美域的谱，只需将 $X$、$Y$ 换成 $X_{\mathrm{red}}$、$Y_{\mathrm{red}}$。一般情形则需缩小 $S$，作有限、根式、满射的基变换 $S'\to S$，再将 $X$、$Y$ 换成其约化概形。有限根式满射不改变平展拓扑，所以这些操作无碍。缩小 $V$ 后可设 $\sF$ 在 $V$ 上局部常值：
\[\xymatrix{V\ar@{^{(}->}[r]^-j&X\ar@{^{(}->}[r]^-f&Y.}\]
用三角
\begin{equation*}\tag{1}\label{VII:eq:2-7-1}
 \cdots\longrightarrow\sF\longrightarrow\eR j_\ast j^\ast\sF\longrightarrow\Delta\longrightarrow\cdots
\end{equation*}
定义 $\Delta$。其上同调层支撑于 $X\setminus V$，且 $\dim(X\setminus V)_\eta<n$，所以由归纳假设可设 $\eR f_\ast\Delta$ 可构造。对\eqref{VII:eq:2-7-1}施以 $\eR f_\ast$，得三角
\[
 \cdots\longrightarrow\eR f_\ast\sF\longrightarrow\eR(fj)_\ast j^\ast\sF\longrightarrow\eR f_\ast\Delta\longrightarrow\cdots.
\]
其中两项可构造且其形成与 $S$ 上任意基变换相容，第三项因而也有此性质。
\subsection{}\label{VII:2-8}
证明\ref{VII:1-9}。问题对 $Y$ 是局部的，可设其仿射。对 $X$ 取仿射覆盖并使用 Leray 谱序列，可进一步设 $X$ 仿射。这些归约保证可以把 $f$ 分解为开嵌入后接固有态射：$f=gj$，于是 $\eR f_\ast=\eR g_\ast\eR j_\ast$。开嵌入由某个\eqref{VII:eq:2-3-1}处理，固有态射由有限性定理处理。
以下推论的证明与第1节相同。
\begin{corollary_}\label{VII:2-9}
在上述定理的假设下，分别对 $K\in\eD_c^b(X,A)$ 或 $K\in\eD_c^b(Y,A)$，存在 $S$ 的非空开集 $U$，使 $\eR f_\ast K$、$\eR f_!K$、$f^\ast K$、$\eR f^!K$ 在 $U$ 上方分别属于 $\eD_c^b(Y,A)$ 或 $\eD_c^b(X,A)$，且其形成与任意基变换 $S'\to U\subset S$ 相容。
\end{corollary_}
\begin{corollary_}\label{VII:2-10}
设 $X$ 如上述定理，且 $A$ 交换。若 $\sF$、$\sG$ 是 $X$ 上的可构造 $A$模层，则在 $S$ 的某个稠密开集 $U$ 上方，各 $\Ext_A^i(\sF,\sG)$ 仍可构造，且其形成与任意基变换 $S'\to U\subset S$ 相容。
\end{corollary_}
\subsection{}\label{VII:2-11}
若 $x$ 为概形 $X$ 的几何点，以 $X_x$ 表示 $X$ 在 $x$ 处的\term{严格Hensel化}{strict henselization}{hensélisé strict}。对 $f:X\to S$ 及 $S$ 的几何点 $t$，以 $X_t$ 表示 $X$ 在 $t$ 处的几何纤维。最后，若 $x$ 为 $X$ 的几何点，$t$ 为 $S_{f(x)}$ 的几何点，则 $(X_x)_t$ 表示 $X_x\to S_{f(x)}$ 在 $t$ 处的纤维。
\begin{definition_}\label{VII:2-12}
设 $f:X\to S$、$K\in\ob\eD^+(X,A)$。若对 $S_{f(x)}$ 的每个几何点 $t$ 都有
\[
 K_x=\eR\Gamma(X_x,K)\iso\eR\Gamma((X_x)_t,K),
\]
则称 $f$\term{在$x$处相对于$K$局部无环}{locally acyclic at $x$ relative to $K$}{localement acyclique en $x$, rel. $K$}。若这对所有 $x$ 都成立，则称 $f$ 相对于 $K$ 局部无环；若任意基变换 $S'\to S$ 后仍成立，则称其相对于 $K$\term{普遍局部无环}{universally locally acyclic}{universellement localement acyclique}。
\end{definition_}
\begin{theorem_}\label{VII:2-13}
设 $f:X\to S$ 为有限型态射，$\sF$ 在 $X$ 上可构造。则存在 $S$ 的稠密开集 $U$，使 $f$ 在 $U$ 上方相对于 $\sF$ 普遍局部无环。
\end{theorem_}
先承认下列结果，其证明见附录。
\begin{lemma_}\label{VII:2-14}
设 $X\xrightarrow fS_1\xrightarrow gS_2$。若 $g$ 光滑，$f$ 相对于 $K$ 普遍局部无环，则 $gf$ 亦然。
\end{lemma_}
归结为 $S$ 整、泛点为 $\eta$ 的情形，对 $\dim X_\eta$ 归纳。首先由归纳假设推出
\begin{equation*}\tag{A}\label{VII:eq:2-14-1}
 \begin{gathered}
 \text{缩小 $S$ 后，存在于 $S$ 上有限的 $T\subset X$，}\\
 \text{使 $f$ 在 $T$ 外普遍局部无环。}
 \end{gathered}
\end{equation*}
这是局部问题，故在 $X$ 上局部化，将 $f$ 分解为 $X\xrightarrow u\dA_S^1\to S$。对 $u$ 用归纳假设，再由\ref{VII:2-14}和通常的论证即可。
证明定理时，因问题是局部的，可设 $X$ 固有。缩小 $S$，使各 $\eR^i f_\ast\sF$ 局部常值且\eqref{VII:eq:2-14-1}适用，然后使用下述引理。
\begin{lemma_}\label{VII:2-15}
设 $f:X\to S$ 固有，$\sF$ 可构造。假设各 $\eR^i f_\ast\sF$ 局部常值，且除某个在 $S$ 上有限的 $T\subset X$ 以外，$f$ 相对于 $\sF$ 局部无环。则 $f$ 相对于 $\sF$ 局部无环。
\end{lemma_}
设 $s$ 为 $S$ 的几何点，$S_s$ 为 $S$ 在 $s$ 处的严格局部化，$t$ 为 $S_s$ 的几何点。设 $X_{(s)}$ 为 $X$ 在 $S_s$ 上的拉回，$\bar j:X_t\to X_{(s)}$、$i:X_s\hookrightarrow X_{(s)}$。仍以 $\sF$ 表示其在 $X_{(s)}$ 或 $X_t$ 上的拉回。须证 $i^\ast\sF\iso i^\ast\eR\bar j_\ast\sF$。令 $\Delta$ 为此映射的映射锥。其上同调层支撑于 $T_s$，故为证明 $\Delta=0$，只需证明 $\eR\Gamma(X_s,\Delta)=0$。后者是下列映射的映射锥：
\begin{align*}
 (\eR f_\ast\sF)_s=\eR\Gamma(X_s,\sF)
 &\longrightarrow\eR\Gamma(X_s,i^\ast\eR\bar j_\ast\sF)\\
 &=\eR\Gamma(X_{(s)},\eR\bar j_\ast\sF)\\
 &=\eR\Gamma(X_t,\sF)=(\eR f_\ast\sF)_t.
\end{align*}
由假设，此特化态射为同构，故 $\Delta=0$。
\begin{corollary_}\label{VII:2-16}
若 $S$ 为域的谱，则任意 $S$概形 $X$ 对每个 $K\in\ob\eD^+(X,A)$ 都普遍局部无环。
\end{corollary_}
由\ref{VII:2-13}取极限即得；这是可行的，因为\ref{VII:2-13}中的 $U$ 此时总等于 $S$。

\origpage{250}{244}
\section{定理\ref{VII:1-1}的证明及附近循环层的可构造性}\label{VII:3}
\subsection{}\label{VII:3-1}
设 $S$ 为\term{严格局部特征线}{strictly local trait}{trait strictement local}，即剩余域可分闭的 Hensel 离散赋值环的谱。记闭点、泛点分别为 $s$、$\eta$，取位于 $\eta$ 上的几何点 $\bar\eta$，即 $k(\eta)$ 的一个可分闭包。设 $X$ 为 $S$概形，$\sF$ 为 $X_\eta$ 上的层。回忆\term{附近循环层}{nearby cycle sheaf}{faisceau de cycles évanescents} $\eR^i\Psi_\eta(\sF)$ 的定义：它是 $X_s$ 上带有 $\gal(\bar\eta/\eta)$ 作用的层。\translatornote{原文将这里的 $\eR\Psi$ 称为 cycles évanescents；依本文给出的定义，它对应现代术语中的“附近循环”。现代通常以 $\eR\Phi$ 表示“消失循环”，后者由从特殊纤维到附近循环的特化态射取锥得到。这里按所定义的函子译名，以免混淆。}
令 $i:X_s\hookrightarrow X$、$j:X_\eta\hookrightarrow X$，并记复合 $X_{\bar\eta}\to X_\eta\hookrightarrow X$ 为 $\bar j$；则
\[
 \eR^i\Psi_\eta(\sF)=i^\ast\eR^i\bar j_\ast(\sF|X_{\bar\eta}).
\]
下面是本节的主要定理，它加强了\cite[XIII 2.3.1、2.4.2]{sga7}。
\begin{theorem_}\label{VII:3-2}
若 $X$ 在 $S$ 上有限型，$\sF$ 可构造，则各附近循环层 $\eR^i\Psi_\eta(\sF)$ 都可构造。
\end{theorem_}
对 $X_\eta$ 的维数 $n$ 归纳。还可设，并且以下总设，$X_\eta$ 在 $X$ 中稠密：将 $X$ 换成 $\overline{X_\eta}$ 不改变各 $\eR^i\Psi$，因为它们支撑于 $\overline{X_\eta}$。
\begin{lemma_}\label{VII:3-3}
若\ref{VII:3-2}在 $\dim X_\eta<n$ 时成立，且 $\dim X_\eta=n$，则存在可构造子层 $\sG_i\subset\eR^i\Psi_\eta(\sF)$，使 $\eR^i\Psi_\eta(\sF)/\sG_i$ 的各局部截面的支撑有限。
\end{lemma_}
设 $s'$ 为 $s$ 上仿射直线的几何泛点，$S'$ 为 $\dA_S^1$ 在 $s'$ 处的严格局部化。它仍是严格局部特征线，且 $S$ 的一致化元也是 $S'$ 的一致化元。
\[\xymatrix{
 &\spec(k')\ar[rr]\ar[dd]\ar[dr]&&(S',\eta',s')\ar[dr]\ar[dd]\\
 \bar\eta'\ar[ur]\ar[dr]&&\dA_{\bar S}^1\ar[d]&&\dA_S^1\ar[dl]\\
 &\bar\eta\ar@{^{(}->}[r]&\bar S\ar[r]&(S,\eta,s)
}\]
令 $k'=k(\bar\eta)\otimes_{k(\eta)}k(\eta')$。若 $\bar S$ 是 $S$ 在 $\bar\eta$ 中的正规化，则 $k'$ 为 $\dA_{\bar S}^1$ 在 $s'$ 处严格局部化的分式域，所以它是域，并有 $\gal(\bar\eta/\eta)=\gal(k'/\eta')$。令 $\bar\eta'$ 为 $k'$ 的一个代数闭包的谱，则 Galois 群 $P=\gal(\bar\eta'/k')$ 是 pro-$p$ 群，其中 $p$ 为 $k(s)$ 的特征指数。
\begin{lemma_}\label{VII:3-4}
对任意 $S'$概形 $X'$ 及 $X'_\eta=X'_{\eta'}$ 上的任意层 $\sF$，$X'/S'$ 和 $X'/S$ 的附近循环层满足
$\eR^i\Psi_\eta(\sF)=\eR^i\Psi_{\eta'}(\sF)^P$。
\end{lemma_}
下图比较了 $X'/S'$ 和 $X'/S$ 的 $\eR\Psi$ 定义中所用的态射：
\[\xymatrix@C=15pt{
 X'_{\bar\eta'}\ar[rr]\ar[d]&&X'_{\eta'}\ar@{=}[dd]\ar@{^{(}->}[dr]&&S'\\
 X'_{\eta'}\times_{\eta'}\spec(k')\ar@{=}[d]\ar[urr]&&&X'\ar[ur]\ar[dr]\\
 X'_{\bar\eta}\ar[rr]&&X'_\eta\ar@{^{(}->}[ur]&&S
}\]
由于 $P$ 是 pro-$p$ 群而 $p$ 在 $A$ 中可逆，Hochschild--Serre 谱序列即给出引理。
现在证明\ref{VII:3-3}。问题是局部的，可设 $X$ 仿射，且 $X\subset\dA_S^N$。设 $f$ 为某个投影 $X\subset\dA_S^N\to\dA_S^1$，$X'=X\times_{\dA_S^1}S'$ 是 $X$ 的“局部化”，$\sF'$ 为 $\sF$ 在 $X'$ 上的拉回：
\[\xymatrix{
 \sF'\text{ 在 }X'\text{ 上}\ar[r]\ar[d]^-\lambda&S'\ar[d]^-\lambda\\
 \sF\text{ 在 }X\text{ 上}\ar[r]^-f&\dA_S^1\ar[r]&S.
}\]
在 $X'_s=X'_{s'}$ 上有
\[
 \lambda^\ast\eR^i\Psi_\eta(\sF)=\eR^i\Psi_\eta(\sF')=\eR^i\Psi_{\eta'}(\sF')^P.
\]
这个层可构造，因为归纳假设适用于 $X'/S'$。再将下述引理用于 $X_s\subset\dA_s^N$ 及附近循环层即可。
\begin{lemma_}\label{VII:3-5}
设 $\dA^N$ 为域 $k$ 上的标准 $N$维仿射空间，$X\subset\dA^N$，$\sF$ 为 $X$ 上的 $A$模层，$\bar\eta$ 为 $\dA^1$ 的几何泛点。记 $X_{\bar\eta,i}$ 为 $X\subset\dA^N\xrightarrow{\pr_i}\dA^1$ 的几何泛纤维，$\sF_{\bar\eta,i}$ 为 $\sF$ 在其上的限制。若所有 $\sF_{\bar\eta,i}$ 可构造，则存在可构造的 $\sF'\subset\sF$，使 $\sF/\sF'$ 的各局部截面支撑有限。
\end{lemma_}
取极限可知，对每个 $i$：
\begin{enumerate}[a)]
\item 存在 $\bar\eta$ 的平展邻域 $U$，及 $X_{U,i}=X\times_{\dA^1,\pr_i}U$ 上的可构造层 $\sH$，使其在 $X_{\bar\eta,i}$ 上的限制为 $\sF_{\bar\eta,i}$。
\item 适当选择 $U$ 后，同构 $\sH_{\bar\eta,i}\iso\sF_{\bar\eta,i}$ 来自 $u:\sH\to\sF_{U,i}$，其中 $\sF_{U,i}$ 为 $\sF$ 在 $X_{U,i}$ 上的拉回。令 $\varphi:X_{U,i}\to X$，则 $u$ 定义 $v_i:\varphi_!\sH\to\sF$，记其像为 $\sF'_i$。有 $(\sF/\sF'_i)_{\bar\eta,i}=0$，取 $\sF'$ 为各 $\sF'_i$ 的和即可。
\end{enumerate}
\subsection{}\label{VII:3-6}
在维数 $n$ 证明\ref{VII:3-2}。问题对 $X_s$ 是局部的，可先设 $X$ 仿射，再取其射影嵌入的闭包而设 $X$ 在 $S$ 上射影。此时有谱序列\cite[I 2.2.3]{sga7}，它是 $\bar j$ 的 Leray 谱序列，再使用固有态射 $X\to S$ 的基变换定理：
\begin{equation*}\tag{1}\label{VII:eq:3-6-1}
 E_2^{pq}=\h^p(X_s,\eR^q\Psi_\eta(\sF))\Rightarrow\h^{p+q}(X_{\bar\eta},\sF).
\end{equation*}
取\ref{VII:3-3}中的 $\sG^q\subset\eR^q\Psi_\eta(\sF)$，记商层为 $\sH^q$。当 $i\ne0$ 时有 $\h^i(X_s,\sH^q)=0$；若 $\h^0(X_s,\sH^q)$ 有限生成，则 $\sH^q$，从而 $\eR^q\Psi_\eta(\sF)$ 都可构造。
现在模去有限生成模来计算\eqref{VII:eq:3-6-1}，即在 $A$模范畴关于有限生成 $A$模构成的厚子范畴的商范畴中计算。由
$0\to\sG^q\to\eR^q\Psi_\eta(\sF)\to\sH^q\to0$ 的上同调长正合列，以及对 $\sG^q$ 应用的\ref{VII:1-10}，得到
$E_2^{pq}\sim\h^p(X_s,\sH^q)$。特别地，$p\ne0$ 时 $E_2^{pq}\sim0$，而
$E_2^{0q}\sim\h^q(X_{\bar\eta},\sF)\sim0$。证明完成。
平行的论证还给出下述结果，它加强了\cite[XIII 2.1.12、2.4.2]{sga7}。
\begin{proposition_}\label{VII:3-7}
附近循环层的形成与特征线的基变换相容。
\end{proposition_}
设 $g:(S',\eta',s')\to(S,\eta,s)$ 为严格局部特征线之间的满射态射，$X/S$，$\sF$ 为 $X_\eta$ 上的挠层，其挠性与剩余特征互素；以 $(X',\sF')$ 表示它们在 $S'$ 上的拉回。与 $g$ 平行的态射，如 $X'_{s'}\to X_s$，仍记为 $g$。须证
\begin{equation*}\tag{3.7.1}\label{VII:eq:3-7-1}
 g^\ast\eR\Psi_\eta(\sF)\iso\eR\Psi_{\eta'}(\sF').
\end{equation*}
该态射由交换图
\[\xymatrix{
 X'_{s'}\ar@{^{(}->}[r]^-{i'}\ar[d]^-g&X'\ar[d]^-g&X'_{\bar\eta'}\ar[l]_-{\bar j'}\ar[d]^-g\\
 X_s\ar@{^{(}->}[r]^-i&X&X_{\bar\eta}\ar[l]_-{\bar j}
}\]
定义。取极限可归结为 $X$ 在 $S$ 上有限型的情形，再对 $\dim X_\eta$ 归纳。问题是局部的，可先设 $X$ 仿射，再设其射影。假设\eqref{VII:eq:3-7-1}在 $\dim X_\eta<n$ 时成立。当 $\dim X_\eta=n$ 时，设 $\Delta$ 为\eqref{VII:eq:3-7-1}的映射锥，则结论由以下两断言推出：
\begin{enumerate}[(A)]
\item $\Delta$ 的上同调层的各局部截面支撑有限。\phantomsection\hypertarget{VII:3-A}{}
\item $\eR\Gamma(X'_{s'},\Delta)=0$，从而由（\hyperlink{VII:3-A}{A}）知 $\Delta=0$。\phantomsection\hypertarget{VII:3-B}{}
\end{enumerate}
\subsection{}\label{VII:3-8}
（\hyperlink{VII:3-A}{A}）的证明与\ref{VII:3-3}平行。在 $X$ 上局部化，考虑映射 $X\to\dA_S^1$；记 $S(x)$ 为 $\dA_S^1$ 在特殊纤维泛点处的严格 Hensel 化，将归纳假设用于 $X_1/S(x)$（其中 $X_1=X\times_{\dA^1}S(x)$）及 $S'(x)$。再取 pro-$p$ 群不变量，便知 $\Delta$ 在 $X'_{s'}\to\dA_s^1$ 的几何泛纤维上为零。此结论在 $X$ 上局部地对每个投影成立，故有（\hyperlink{VII:3-A}{A}）。
\subsection{}\label{VII:3-9}
对（\hyperlink{VII:3-B}{B}），$X$ 在 $S$ 上固有，保证
\[
 \eR\Gamma(X'_{s'},\eR\Psi_{\eta'}(\sF'))=\eR\Gamma(X'_{\bar\eta'},\sF')
\]
和
\[
 \eR\Gamma(X_s,\eR\Psi_\eta(\sF))=\eR\Gamma(X_{\bar\eta},\sF).
\]
由可分闭基域扩张下的不变性，
\[
 \eR\Gamma(X_s,\eR\Psi_\eta(\sF))\iso\eR\Gamma(X'_{s'},g^\ast\eR\Psi_\eta(\sF)).
\]
所以 $\eR\Gamma(X'_{s'},\Delta)$ 是
$\eR\Gamma(X_{\bar\eta},\sF)\to\eR\Gamma(X'_{\bar\eta'},\sF')$ 的映射锥；由同一不变性定理，此映射是同构。
\subsection{}\label{VII:3-10}
证明\ref{VII:1-1}。可设 $S$ 连通，从而整。若 $\dim S=0$，即 $S$ 为域的谱，应用\ref{VII:1-9}、\ref{VII:1-10}即可。若 $\dim S=1$，定理\ref{VII:1-9}保证各 $\eR^i f_\ast\sF$ 在 $S$ 的有限点集 $T$ 的补集上方可构造。只剩证明其在各 $Y_t$（$t\in T$）上的限制可构造。在 $t$ 处局部化后证明即可，所以可设 $S$ 为严格局部特征线。关键情形如下。
\begin{lemma_}\label{VII:3-11}
设 $X$ 在 $S$ 上有限型，$j:X_\eta\hookrightarrow X$，$i:X_s\hookrightarrow X$，$\sF$ 在 $X_\eta$ 上可构造。则 $i^\ast\eR j_\ast\sF$ 可构造。
\end{lemma_}
令 $I=\gal(\bar\eta/\eta)$ 为\term{惯性群}{inertia group}{groupe d'inertie}。已知它是 pro-$p$ 群 $P$ 对 $\widehat\dZ_{p'}(1)$ 的扩张，其中 $p$ 为剩余域的特征指数。Hochschild--Serre 谱序列给出
\[
 E_2^{pq}=\sH^p(I,\eR^q\Psi_\eta(\sF))\Rightarrow i^\ast\eR^{p+q}j_\ast\sF.
\]
令 $R_t^q=\eR^q\Psi_\eta(\sF)^P$，下标 $t$ 表示\term{温和}{tame}{modéré}。若 $\sigma$ 为 $\widehat\dZ_{p'}(1)$ 的一个拓扑生成元，则
\[
 E_2^{0q}=\ker(\sigma-1,R_t^q),\qquad
 E_2^{1q}=\operatorname{coker}(\sigma-1,R_t^q),\qquad
 E_2^{pq}=0\quad(p\ne0,1).
\]
附近循环层可构造，所以这些层也可构造，引理\ref{VII:3-11}得证。
\subsection{}\label{VII:3-12}
现在处理一般情形，仍设 $S$ 为严格 Hensel 特征线。若 $X=X_\eta$，则 $f$ 为 $X\to Y_\eta\to Y$ 的复合，应用\ref{VII:1-9}及\ref{VII:3-11}即可。一般地，取 $j:X_\eta\hookrightarrow X$，令 $\Delta$ 为 $\sF\to\eR j_\ast j^\ast\sF$ 的映射锥。其上同调层支撑于 $X_s$；将\ref{VII:1-9}用于 $X_s\to Y_s$，得 $\eR f_\ast\Delta$ 可构造。复形 $\eR f_\ast\eR j_\ast j^\ast\sF=\eR(fj)_\ast j^\ast\sF$ 也可构造；再由三角
\[
 \cdots\longrightarrow\eR f_\ast\sF\longrightarrow\eR f_\ast\eR j_\ast j^\ast\sF\longrightarrow\eR f_\ast\Delta\longrightarrow\cdots
\]
得到结论。

\origpage{256}{250}
\section{局部双重对偶}\label{VII:4}
\subsection{}\label{VII:4-1}
设 $S$ 正则、维数为 $0$ 或 $1$，令 $K_S=A$，即常值为 $A$ 的常层。对 $K\in\ob\eD_{\mathrm{ctf}}^b(S,A)$，令
$DK=\rHom(K,K_S)\in\ob\eD^+(S,A^\circ)$，这是右 $A$模层复形。仍有 $DK\in\ob\eD_{\mathrm{ctf}}^b(S,A^\circ)$；由于 $S$ 的维数不超过 $1$，可以显式计算得到 $K\iso DDK$。当 $A=\dZ/n$ 时，见\hyperref[V]{《对偶》}\ref{V:1-4}。
\subsection{}\label{VII:4-2}
以下设 $A$ 交换，这个假设想必并非必要。对有限型态射 $a:X\to S$，令 $K_X=\eR a^!K_S$。对 $K\in\ob\eD_{\mathrm{ctf}}^b(X,A)$，令
$DK=\rHom(K,K_X)\in\ob\eD_{\mathrm{ctf}}^b(X,A)$。我们的主要结果是：
\begin{theorem_}\label{VII:4-3}
有 $K\iso DDK$，其中 $K\in\ob\eD_{\mathrm{ctf}}^b(X,A)$。
\end{theorem_}
\begin{lemma_}\label{VII:4-4}
若 $a$ 固有，则 $\eR a_\ast K\iso\eR a_\ast DDK$。
\end{lemma_}
Poincaré 对偶
\[
 \eR a_\ast\rHom(K,\eR a^!K_S)=\rHom(\eR a_!K,K_S)
\]
给出同构 $\eR a_\ast D=D\eR a_!=D\eR a_\ast$。引理由下图的交换性得到，参见\hyperref[V]{《对偶》}\ref{V:1-2}：
\[\xymatrix{
 \eR a_\ast K\ar[r]\ar[d]^-\wr&\eR a_\ast DDK\ar[d]^-\wr\\
 DD\eR a_\ast K\ar[r]^-\sim&D\eR a_\ast DK.
}\]
\subsection{}\label{VII:4-5}
局部化后可知，只需处理 $S$ 为域的谱或特征线的情形。若 $S$ 为域的谱，对 $\dim X$ 归纳。问题是局部的，可设 $X$ 固有。此外，由归纳假设及通常的论证，$K\to DDK$ 的映射锥 $\Delta$ 的上同调层都是\term{摩天楼层}{skyscraper sheaf}{faisceau en gratte-ciel}。引理保证 $\eR a_\ast\Delta=0$，从而 $\Delta=0$。
\subsection{}\label{VII:4-6}
若 $S$ 为特征线，已由\ref{VII:4-5}知道\ref{VII:4-3}在泛纤维上成立。也可设 $X$ 在 $S$ 上固有，并对 $\dim X_s$ 归纳。映射锥 $\Delta$ 的上同调层仍是 $X_s$ 上的摩天楼层，因而与上面一样得到结论。
\subsection{}\label{VII:4-7}
设 $A=\dZ/m$。上述理论此时还有一种变体：可取任意 $K\in\eD_c^b(X,A)$。复形 $\eR a^!\dZ/m$ 的内射维数有限，见\cite[XVIII 3.1.7]{sga4}，所以 $DK$ 仍在 $\eD_c^b(X,A)$ 中。证明 $S$ 上的局部双重对偶时，使用 $\dZ/m$ 是 $\dZ/m$模的对偶化模这一事实；之后与前面同样进行。

\origpage{258}{252}
\section*{附录\quad L. Illusie 撰}\label{VII:5}
\addcontentsline{toc}{section}{附录：L. Illusie}
本附录采用了原 SGA 5 II 的若干部分及 P. Deligne 致作者的一封信，证明本篇\ref{VII:2-14}，以及特化定理\cite[XVI 2.1]{sga4}和\hyperref[I]{《平展上同调》}\ref{I:5-1-7}的若干推广与变体。
以下所考虑的概形及态射均假设拟紧且拟分离。若 $X$ 为概形，$x$ 为其几何点，以 $X(x)$ 表示 $X$ 在 $x$ 处的严格局部化。
固定基概形 $S$ 及其上的环层 $A$，不要求 $A$ 交换或诺特。若 $X$ 为 $S$概形，将 $\eD(X,A_X)$ 简记为 $\eD(X)$。
\setcounter{section}{0}
\section{上同调固有性}\label{VII:5-1}
\begin{proposition_}\label{VII:5-1-1}
设 $f:X\to Y$ 为 $S$态射，$E\in\ob\eD^+(X)$。以下条件等价：
\begin{enumerate}[\indent (i)]
\item $\eR f_\ast E$ 的形成与任意有限基变换 $S'\to S$ 交换。
\item $\eR f_\ast E$ 的形成与任意拟有限基变换 $S'\to S$，以及拟有限态射的逆极限所给出的基变换交换。
\item 对每个\term{特化态射}{specialization morphism}{flèche de spécialisation} $i:t\to S(s)$（\hyperref[I]{《平展上同调》}\ref{I:5-1-2}），令 $\bar S$ 为 $S(s)$ 中 $i(t)$ 的闭包这一整概形在 $k(t)$ 中的正规化。根据\cite[IV 6.15.5、18.5.11、18.8.16]{ega}，$\bar S$ 是局部整概形，其闭点在 $s$ 上为根式的。设 $\bar f:\bar X\to\bar Y$ 为沿 $\bar S\to S$ 对 $f$ 作基变换所得的态射，$f_s:X_s\to Y_s$ 为 $f$ 在 $s$ 处的纤维；则基变换态射
\[
 (\eR\bar f_\ast(E|\bar X))|Y_s\longrightarrow\eR{f_s}_\ast(E|X_s)
\]
为同构。
\end{enumerate}
\end{proposition_}
由 Zariski 主定理知 (i) 与 (ii) 等价；(ii) 显然蕴含 (iii)。证明 (iii) 蕴含 (ii)。对任意作为拟有限态射逆极限的 $S'\to S$ 及 $S'$ 的任意几何点 $s$，以 $C(S',s)$ 表示基变换态射
\[
 (\eR{f'}_\ast(E|X'))|Y_s\longrightarrow\eR{f_s}_\ast(E|X_s)
\]
的\term{映射锥}{mapping cone}{cône}，其中 $f':X'\to Y'$ 由 $S'\to S$ 对 $f$ 作基变换得到，$f_s:X_s\to Y_s$ 为在 $s$ 处的纤维。须证每个这样的 $(S',s)$ 都有 $C(S',s)$ 无环。对 $N$ 归纳证明 $\h^iC(S',s)=0$（$i\leqslant N$）。由主定理、取极限及去除根式态射，可设 $S$、$S'$ 都严格局部，$S'$ 整且闭点为 $s$，$S'\to S$ 有限。取 $S'$ 的几何泛点 $t$，令 $\bar S$ 为 $S'$ 在 $k(t)$ 中的正规化。令 $\bar S_\bullet$ 为 $\bar S\to S'$ 的余骨架单纯概形：
\[
 \bar S_0=\bar S,\qquad \bar S_n=(\bar S/S')^{n+1},\qquad\ldots,
\]
并记 $\bar X_\bullet=X\times_S\bar S_\bullet$，$g_\bullet:\bar X_\bullet\to X'$ 为典范投影。对用于表示 $E$ 的层或复形 $G$，记 $G|X'=G'$。双复形 ${g_\bullet}_\ast g_\bullet^\ast G'$ 的各列为 ${g_n}_\ast g_n^\ast G'$，它给出 $G'$ 的一个分解，见\cite[VIII 8]{sga4}。所以可用各列等同于 $C(\bar S_n,s)$ 的双复形表示 $C(S',s)$。由 (iii)，$C(\bar S_0,s)=C(\bar S,s)$ 无环；当 $n>0$ 时，由归纳假设有 $\h^iC(\bar S_n,s)=0$（$i\leqslant N$）。因此 $\h^iC(S',s)=0$ 对 $i\leqslant N+1$ 成立，证明完成。
\subsection{注}\label{VII:5-1-2}
设 $(f,E)$ 满足\ref{VII:5-1-1}的条件 (i)。
\begin{enumerate}[a)]
\item 若 $E\in\ob\eD^b(X)$ 且 $f_\ast$ 的上同调维数有限，则对任意 $F\in\ob\eD^-(S,A^\circ)$，典范态射
\begin{equation*}\tag{$*$}\label{VII:eq:5-1-1-1}
 \eR f_\ast(E)\lotimes_A q^\ast F\longrightarrow\eR f_\ast(E\lotimes_A p^\ast F)
\end{equation*}
是同构，其中 $p:X\to S$、$q:Y\to S$ 为投影。对 $E\in\ob\eD^-(X)$ 也有同样结论。
\end{enumerate}
事实上，分解归约后，只需对 $F=i_!A_{S'}$ 验证断言，其中 $i:S'\to S$ 拟有限。由主定理可设 $i$ 有限，此时结论等价于 $\eR f_\ast(E)$ 与沿 $i$ 的基变换交换。
\begin{enumerate}[a)]\setcounter{enumi}{1}
\item 设 $A$ 诺特且可构造。则对每个可构造且 Tor 维数有限的 $F\in\ob\eD^b(S,A^\circ)$，上面的态射\eqref{VII:eq:5-1-1-1}都是同构。
\end{enumerate}
事实上，分解归约后可设 $F=i_!M$，其中 $i:S'\to S$ 局部闭，$A$ 在 $S'$ 上局部常值，$M$ 可构造、Tor 维数有限，且其上同调在 $S'$ 上局部常值。将基从 $S$ 换到 $S'$，便归结为 $F$\term{完美}{perfect}{parfait}的情形，参见\cite[I]{sga6}；再经局部化及分解归约，最终归结为 $F=A$，此时结论显然。
\subsection{}\label{VII:5-1-3}
设 $f:X\to Y$ 为 $S$态射，$E\in\ob\eD^+(X)$。若 $\eR f_\ast E$ 的形成与每个基变换 $S'\to S$ 交换，就称 $(f,E)$ 相对于 $S$\term{上同调固有}{cohomologically proper}{cohomologiquement propre}。等价地，对任意基变换 $S'\to S$ 后所得的 $f'=f\times_SS'$ 及 $E'=E|X\times_SS'$，$\eR{f'}_\ast E'$ 的形成都与任意有限基变换 $S''\to S'$ 交换，参见\cite[XII 6.1]{sga4}。下面给出几个上同调固有性的例子；集合层或非交换群层的类似概念见\cite[XIII]{sga1}。
\paragraph*{1.3.1.}
设 $A$ 为挠环层，$f$ 固有。则对每个 $E\in\ob\eD^+(X)$，$(f,E)$ 相对于 $S$ 上同调固有，这是固有基变换定理\cite[XII 5.1]{sga4}。
\paragraph*{1.3.2.}
设 $S$ 为有限概形，$X$、$Y$ 在 $S$ 上有限型，且 $A$ 被某个在 $S$ 上可逆的整数零化。则对每个 $E\in\ob\eD^+(X)$，$(f,E)$ 相对于 $S$ 上同调固有，见本篇\ref{VII:1-9}。
\paragraph*{1.3.3.}
设 $A$ 为常值诺特环，被某个在 $S$ 上可逆的整数 $n$ 零化；设 $D$ 为 $Y$ 上相对于 $S$ 具有\term{正规交叉}{normal crossings}{croisements normaux}的除子，$f:X=Y\setminus D\to Y$ 为其补集的包含映射，参见\cite[XIII 2.1]{sga1}。设 $E$ 为 $X$ 上的 $A$模层，满足以下条件之一：
\begin{enumerate}[(i)]
\item $E$ 局部常值可构造，且沿 $D$ 温和分歧；
\item 对 $Y$ 及 $S$ 上的平展拓扑，$E$ 局部地来自 $S$ 上某个 $A$模层的拉回。
\end{enumerate}
则 $(f,E)$ 相对于 $S$ 上同调固有，且 $\eR f_\ast E$ 可构造；可与\cite[XIII 2.4]{sga1}比较。\translatornote{后一项可构造性在 (ii) 中还需有限性假设。例如取 $S=\operatorname{Spec}k$，$k$ 为代数闭域；取结构态射 $q:Y=\mathbf A^1_k\to S$，$D=\{0\}$，$f:X=\mathbf G_{m,k}\hookrightarrow Y$，$p=qf$，以及 $A=\mathbf Z/n$，其中 $n>1$ 且与 $\operatorname{char}k$ 互素。令 $M=\bigoplus_{r\geqslant1}A$，取 $E=p^*\underline M$，则满足 (ii)，但 $f_*E$ 在 $X$ 上的茎为无限生成的 $M$，故非可构造。上同调固有性不需此有限性条件：由下面的纯性公式，$\eR f_*(p^*M)=\eR f_*(\mathbf Z/n)\mathbin{\overset{\mathsf L}{\otimes}}q^*M$，对任意基变换同样成立。若另假设 (ii) 中的 $E$ 可构造，例如局部拉回的 $M$ 可构造，则按 $M$ 的有限分层归约到有限生成常值模，所列纯性公式也给出所需可构造性。}
简述 (i) 的证明。问题局部于 $Y$ 的一个点 $y$ 的邻域，可设 $D$ 为光滑除子之和 $D=\sum_{i=1}^mD_i$。由相对 Abhyankar 引理\cite[XIII 5.5]{sga1}，可在 $y$ 附近用形如
\[
 \widetilde Y=Y[T_1,\dots,T_r]/(T_1^{n_1}-t_1,\dots,T_r^{n_r}-t_r)
\]
的有限覆盖 $\widetilde Y\to Y$ 平凡化 $E$；其中 $t_i$ 为通过 $y$ 的各光滑除子的局部方程，$n_i$ 与 $k(y)$ 的特征互素。$D$ 在 $\widetilde Y$ 中的拉回 $\widetilde D$ 仍相对于 $S$ 正规交叉；若 $g:\widetilde X=\widetilde Y\setminus\widetilde D\to X$ 为投影，则 $g^\ast E$ 常值。因为 $E$ 嵌入 $g_\ast g^\ast E$，且商仍温和分歧，简单的分解归约便将问题归结为 $E$ 常值的情形。
记投影为 $p:X\to S$、$q:Y\to S$，对 $1\leqslant i\leqslant m$，记典范包含映射为 $f_i:Y\setminus D_i\to Y$。\translatornote{原书此处将 $f_i$ 的目标排作 $D$，应为 $Y$；否则“包含映射”及随后的 $Y$ 上的直像层均无从定义。}
由相对纯性定理\cite[XVI 3.7]{sga4}，对 $m$ 归纳容易推出：对 $S$ 上任意局部常值 $A$模层 $M$，典范态射
\begin{equation*}\tag{1.3.3.1}\label{VII:eq:5-1-3-3-1}
 \eR f_\ast(\dZ/n)\lotimes q^\ast M\longrightarrow\eR f_\ast(p^\ast M)
\end{equation*}
是同构；另一方面有
\begin{equation*}\tag{1.3.3.2}\label{VII:eq:5-1-3-3-2}
 \begin{aligned}
 {f_i}_\ast(\dZ/n)&=\dZ/n,\\
 \eR^1{f_i}_\ast(\dZ/n)&=(\dZ/n)(-1)_{D_i}.
 \end{aligned}
\end{equation*}
第二个典范同构由基本类 $\cl(D_i)$ 给出，参见\hyperref[IV]{《循环》}\ref{IV:2-1-4}；此外，
\begin{align*}
 \eR^q{f_i}_\ast(\dZ/n)&=0\quad(q>1),\\
 \eR f_\ast(\dZ/n)&=\bigotimes_i^{\mathsf L}\eR{f_i}_\ast(\dZ/n).
\end{align*}
也就是说，
\[
 \eR^1f_\ast(\dZ/n)=\bigoplus_{i=1}^m(\dZ/n)(-1)_{D_i},\qquad
 \textstyle\bigwedge^\bullet\eR^1f_\ast(\dZ/n)\iso\eR^\bullet f_\ast(\dZ/n).
\]
所以在 (i) 时，1.3.3 的结论由\eqref{VII:eq:5-1-3-3-1}及\eqref{VII:eq:5-1-3-3-2}推出。在 (ii) 时，取极限可归结为 $E=p^\ast M$、$M$ 可构造的情形；再通过分解归约和使用 (i) 的基变换，归结为已经处理的 $M$ 常值情形。
由上述证明还知，对任意 $M\in\ob\eD^+(S)$，\eqref{VII:eq:5-1-3-3-1}都是同构。
可将\eqref{VII:eq:5-1-3-3-2}与 $S=\spec(\dC)$ 时整系数上同调中的类似公式比较，例如见\cite[3.1]{de71-2}。那些公式来自如下事实：在经典拓扑下局部地，$D$ 的补集与一个环面同伦等价。
取极限可从\eqref{VII:eq:5-1-3-3-2}得到 $\ell$进上同调的类似公式，即将 $\dZ/n$ 换成 $\dZ_\ell$，其中素数 $\ell$ 在 $S$ 上可逆。设 $q:Y\to S$ 固有且光滑，则由 Weil 猜想\cite{de74,de80}可知，Leray 谱序列
\[
 E_2^{ij}=\eR^i q_\ast\eR^j f_\ast\dZ_\ell\Rightarrow\eR^\bullet p_\ast\dZ_\ell
\]
模去挠性后在 $E_3$ 页退化，见\cite[第6节]{de71-1}。当 $S=\spec(\dC)$ 时，整系数上同调中也有同样现象，参见同处及\cite[3.2.13]{de71-2}。

\origpage{262}{256}
\section{特化与余特化}\label{VII:5-2}
设 $f:X\to S$，$E\in\ob\eD^+(X)$。
\subsection{}\label{VII:5-2-1}
假设 $\eR f_\ast E$ 的形成与任意有限基变换交换。对 $S$ 的几何点之间的每个特化 $u:t\to S(s)$，可如下定义\term{特化态射}{specialization morphism}{flèche de spécialisation}
\begin{equation*}\tag{2.1.1}\label{VII:eq:5-2-1-1}
 \spe(u):\eR\Gamma(X_s,E|X_s)\longrightarrow\eR\Gamma(X_t,E|X_t).
\end{equation*}
如\ref{VII:5-1-1}(iii)，令 $\bar S$ 为 $S(s)$ 中 $u(t)$ 的闭包在 $k(t)$ 中的正规化，$\bar f:\bar X\to\bar S$ 为沿 $\bar S\to S$ 对 $f$ 作基变换所得的态射。由\ref{VII:5-1-1}中显然的方向 (i)$\Rightarrow$(iii)，基变换态射
\begin{equation*}\tag{2.1.2}\label{VII:eq:5-2-1-2}
 \eR\Gamma(\bar X,E|\bar X)=(\eR\bar f_\ast(E|\bar X))_s\longrightarrow\eR\Gamma(X_s,E|X_s)
\end{equation*}
是同构。以\eqref{VII:eq:5-2-1-2}的逆与限制态射
\begin{equation*}\tag{2.1.3}\label{VII:eq:5-2-1-3}
 \eR\Gamma(\bar X,E|\bar X)\longrightarrow\eR\Gamma(X_t,E|X_t)
\end{equation*}
的复合定义\eqref{VII:eq:5-2-1-1}。
\subsection{}\label{VII:5-2-2}
设 $f$ 相对于 $E$ 局部无环，见本篇\ref{VII:2-12}。对 $S$ 的几何点之间的每个特化 $u:t\to S(s)$，如下定义\term{余特化态射}{cospecialization morphism}{flèche de cospécialisation}
\begin{equation*}\tag{2.2.1}\label{VII:eq:5-2-2-1}
 \cosp(u):\eR\Gamma(X_t,E|X_t)\longrightarrow\eR\Gamma(X_s,E|X_s).
\end{equation*}
设 $g:X_t\to X$、$g':\bar X\to X$ 为典范态射。由局部无环性假设，典范态射
\begin{equation*}\tag{2.2.2}\label{VII:eq:5-2-2-2}
 \eR{g'}_\ast(E|\bar X)\longrightarrow\eR g_\ast(E|X_t)
\end{equation*}
是同构。事实上，可设 $S$ 严格局部且整，闭点为 $s$，$t$ 为几何泛点。计算\eqref{VII:eq:5-2-2-2}在 $X$ 的几何点 $x$ 处的茎，设 $x$ 在 $S$ 中的像为 $y$。将 $S$ 换成在 $y$ 处的严格局部化后，可设 $y=s$。因为 $\bar S\to S$ 是有限态射的逆极限，且在 $s$ 处根式，可得 $(\eR{g'}_\ast(E|\bar X))_x=E_x$，并且\eqref{VII:eq:5-2-2-2}在 $x$ 处的茎等同于限制态射
\[
 E_x=\eR\Gamma(X(x),E)\longrightarrow\eR\Gamma(X(x)_t,E|X(x)_t),
\]
由假设它是同构。对\eqref{VII:eq:5-2-2-2}施以 $\eR\Gamma(X,-)$，得限制态射\eqref{VII:eq:5-2-1-3}为同构。以\eqref{VII:eq:5-2-1-3}的逆与基变换态射\eqref{VII:eq:5-2-1-2}的复合定义\eqref{VII:eq:5-2-2-1}。
\subsection{}\label{VII:5-2-3}
若 $s''\xrightarrow{u'}s'\xrightarrow us$ 为 $S$ 的几何点之间的特化态射，则只要两边有定义，便可验证
\[
 \spe(u')\spe(u)=\spe(uu'),\qquad
 \cosp(u)\cosp(u')=\cosp(uu').
\]
\subsection{}\label{VII:5-2-4}
设 $\eR f_\ast E$ 的形成与任意有限基变换交换，且 $f$ 相对于 $E$ 局部无环。由定义可知，对 $S$ 的几何点之间的每个特化 $u:t\to S(s)$，$\spe(u)$ 和 $\cosp(u)$ 都是同构，且互为逆。
例如，当 $A$ 诺特且被某个在 $S$ 上可逆的整数 $n$ 零化，$f$ 固有且光滑，$E$ 的上同调局部常值时，上述假设就成立：使用固有态射的上同调固有性（\ref{VII:5-1-3}）及光滑态射的局部无环性\cite[XV]{sga4}。这样便恢复了“特化定理”\cite[XVI 2.1]{sga4}。
另一方面，\ref{VII:5-1-1}有如下推论。
\begin{proposition_}\label{VII:5-2-5}
设 $f$ 相对于 $E$ 局部无环，且对每个特化 $u:t\to S(s)$，余特化态射 $\cosp(u)$ 都是同构。则 $\eR f_\ast E$ 的形成与任意有限基变换交换。
\end{proposition_}
结合\ref{VII:5-2-4}，可知\ref{VII:5-2-5}推广了\hyperref[I]{《平展上同调》}\ref{I:5-1-7}。
\begin{corollary_}\label{VII:5-2-6}
设 $f$ 相对于 $E$ 局部无环。令 $x$ 为 $X$ 的几何点，像为 $S$ 中的 $s$，以 $f_{(x)}:X(x)\to S(s)$ 表示 $f$ 诱导的态射。则 $\eR{f_{(x)}}_\ast(E|X(x))$ 的形成与任意有限基变换 $S'\to S(s)$ 交换。
\end{corollary_}
根据\ref{VII:5-2-5}，只需验证 $(f_{(x)},E|X(x))$ 的余特化态射是同构。由余特化态射的传递性（\ref{VII:5-2-3}），只需证明：对特化 $u:t\to S(s)$，
\[
 \cosp(u):\eR\Gamma(X(x)_t,E|X(x)_t)\longrightarrow\eR\Gamma(X(x)_s,E|X(x)_s)
\]
是同构。因为 $s$ 在 $S(s)$ 中闭，$X(x)_s$ 是以 $x$ 为闭点的严格局部概形，故 $\eR\Gamma(X(x)_s,E|X(x)_s)=E_x$；$\cosp(u)$ 为同构，它是限制同构 $E_x\to\eR\Gamma(X(x)_t,E|X(x)_t)$ 的逆。
\begin{corollary_}\label{VII:5-2-7}
设 $f:X\to Y$ 为 $S$态射，$p:X\to S$、$q:Y\to S$ 为投影，$E\in\ob\eD^+(X)$。假设 $A$ 局部常值，$f$ 相对于 $E$ 局部无环，且对 $Y$ 的每个几何点 $y$，$q$ 在 $y$ 处相对于 $y$ 附近的任意局部常值层都局部无环。则 $p$ 相对于 $E$ 局部无环。
\end{corollary_}
设 $x$ 为 $X$ 的几何点，在 $Y$、$S$ 中的像分别为 $y$、$s$，$t\to S(s)$ 为特化，以 $f_{(x)}:X(x)\to Y(y)$ 表示 $f$ 诱导的态射。须证限制态射
\begin{equation*}\tag{$*$}\label{VII:eq:5-2-7-1}
 E_x\longrightarrow\eR\Gamma(X(x)_t,E|X(x)_t)
\end{equation*}
是同构。可设 $A$ 常值。由\ref{VII:5-2-4}及\ref{VII:5-2-6}，对每个特化 $z\to Y(y)$，特化态射
\begin{align*}
 E_x=(\eR{f_{(x)}}_\ast(E|X(x)))_y
 &\longrightarrow(\eR{f_{(x)}}_\ast(E|X(x)))_z\\
 &=\eR\Gamma(X(x)_z,E|X(x)_z)
\end{align*}
是同构；换言之，$\eR{f_{(x)}}_\ast(E|X(x))$ 常值，值为 $E_x$。对 $q$ 的无环性假设因而说明限制态射
\begin{align*}
 E_x&\longrightarrow\eR\Gamma(Y(y)_t,(E_x)|Y(y)_t)\\
 &=\eR\Gamma\bigl(Y(y)_t,(\eR{f_{(x)}}_\ast(E|X(x)))|Y(y)_t\bigr)
\end{align*}
是同构。而该态射等同于\eqref{VII:eq:5-2-7-1}，所以结论成立。
特别地，当 $q$ 光滑，且 $A$ 被某个在 $S$ 上可逆的整数 $n$ 零化时，\ref{VII:5-2-7}中对 $q$ 的假设成立，见\cite[XV]{sga4}。另一方面，如果\ref{VII:5-2-7}的假设在任意基变换 $S'\to S$ 后都成立，就得到 $p$ 相对于 $E$ 普遍局部无环。这些观察特别给出了本篇\ref{VII:2-14}。
\begin{corollary_}\label{VII:5-2-8}
设 $f:X\to S$，$E\in\ob\eD^+(X)$。假设 $f$ 相对于 $E$ 普遍局部无环，$\eR f_\ast E$ 的形成与任意有限基变换交换；并且对每个态射 $s'\to s$，其中 $s$ 是 $S$ 的几何点、$s'$ 是可分闭域的谱，基变换态射
$\eR\Gamma(X_s,E|X_s)\to\eR\Gamma(X_{s'},E|X_{s'})$ 都是同构。则 $(f,E)$ 上同调固有，见\ref{VII:5-1-3}。
\end{corollary_}
设 $f':X'\to S'$ 为沿基变换 $S'\to S$ 从 $f$ 得到的态射。须证 $\eR{f'}_\ast(E|X')$ 的形成与任意有限基变换交换。根据\ref{VII:5-2-5}，只需证明：对每个特化 $u':t'\to S'(s')$，余特化态射
\[
 \cosp(u'):\eR\Gamma(X'_{t'},E|X'_{t'})\longrightarrow\eR\Gamma(X'_{s'},E|X'_{s'})
\]
是同构。设 $s$、$t$ 分别为 $s'$、$t'$ 在 $S$ 中的像几何点，$u:t\to S(s)$ 为 $u'$ 确定的特化。容易看出有交换方块
\[\xymatrix{
 \eR\Gamma(X_t,E|X_t)\ar[r]^-{\cosp(u)}\ar[d]&\eR\Gamma(X_s,E|X_s)\ar[d]\\
 \eR\Gamma(X'_{t'},E|X'_{t'})\ar[r]^-{\cosp(u')}&\eR\Gamma(X'_{s'},E|X'_{s'})
}\]
其中竖箭头为基变换态射，按假设都是同构；又由\ref{VII:5-2-4}，$\cosp(u)$ 是同构，故 $\cosp(u')$ 也是同构，证明完成。
下面是\ref{VII:5-2-8}的一个应用。设 $A$ 被某个在 $S$ 上可逆的整数 $n$ 零化，$f$ 局部有限型且相对于 $E$ 普遍局部无环。设 $x$ 为 $X$ 的几何点，像为 $S$ 中的 $s$，$f_{(x)}:X(x)\to S(s)$ 为 $f$ 诱导的态射。由\ref{VII:5-2-6}及本篇\ref{VII:1-9}，通过取极限可知，$(f_{(x)},E|X(x))$ 满足\ref{VII:5-2-8}的假设，故上同调固有。
\subsection{}\label{VII:5-2-9}
设 $A$ 常值，$f:X\to S$，$E\in\ob\eD^+(X)$ 为 Tor 维数有限的复形。若对 $X$ 的每个几何点 $x$，记其在 $S$ 中的像为 $s$，对每个特化 $t\to S(s)$ 及每个 $A^\circ$模 $M$，限制态射
\[
 E_x\lotimes_A M\longrightarrow\eR\Gamma\bigl(X(x)_t,(E|X(x)_t)\lotimes_A M\bigr)
\]
都是同构，就称 $f$ 相对于 $E$\term{强局部无环}{strongly locally acyclic}{fortement localement acyclique}。我不知道“局部无环”是否蕴含“强局部无环”。
\begin{proposition_}\label{VII:5-2-10}
在\ref{VII:5-2-9}的假设下，另设 $A$ 为诺特挠环，$S$ 诺特，且 $f$ 相对于 $E$ 强局部无环。则对任意各方块为笛卡尔方块的图
\[\xymatrix{
 X\ar[d]^-f&X'\ar[l]_-j\ar[d]^-{f'}&X''\ar[l]_-{j'}\ar[d]^-{f''}\\
 S&S'\ar[l]_-i&S''\ar[l]_-{i'},
}\]
其中 $i$ 拟有限、$i'$ 为开浸入，以及任意 $F\in\ob\eD^+(S'',A^\circ)$，典范态射
\begin{equation*}\tag{$*$}\label{VII:eq:5-2-10-1}
 j^\ast E\lotimes {f'}^\ast\eR{i'}_\ast F\longrightarrow\eR{j'}_\ast\bigl((jj')^\ast E\lotimes {f''}^\ast F\bigr)
\end{equation*}
为同构。
\end{proposition_}
可将此陈述与\cite[XV 1.17]{sga4}比较。
证明\ref{VII:5-2-10}。容易归结为 $S=S'$ 严格局部、$F$ 可构造且集中在 $0$ 次的情形。然后用形如 $u_\ast M$ 的层的有限直和给出 $F$ 的右分解，其中 $u:t\to S$ 为几何点，$M$ 为 $A^\circ$模。由假设，\eqref{VII:eq:5-2-10-1}对 $F=u_\ast M$ 是同构，结果即得。
\origpage{267}{261}
\subsection*{参考文献}
\begin{description}
\item[{[1]}] P. Deligne, \emph{Théorie de Hodge I}（Hodge 理论 I）, Actes Cong. int. math., 1970, t. I, 425--430.
\item[{[2]}] P. Deligne, \emph{Théorie de Hodge II}（Hodge 理论 II）, Pub. Math. IHES n\textsuperscript{o} 40, 1972, 1--57.
\item[{[3]}] P. Deligne, \emph{La conjecture de Weil I}（Weil 猜想 I）, Pub. Math. IHES n\textsuperscript{o} 43, 1974, 273--307.
\item[{[4]}] P. Deligne, \emph{La conjecture de Weil II}（Weil 猜想 II）, 待刊。
\end{description}
